\documentclass[12pt,a4paper,oneside]{report}

\usepackage[T1]{fontenc}
\usepackage[utf8]{inputenc}
\usepackage{lmodern}
\usepackage{geometry}
\usepackage{setspace}
\usepackage{amsmath,amssymb}
\usepackage{booktabs}
\usepackage{array}
\usepackage{longtable}
\usepackage{multirow}
\usepackage{graphicx}
\usepackage{url}
\usepackage{hyperref}
\usepackage[numbers,sort&compress]{natbib}

\geometry{margin=1in}
\onehalfspacing
\hypersetup{
  colorlinks=true,
  linkcolor=black,
  citecolor=black,
  urlcolor=blue
}

\begin{document}

\begin{titlepage}
\centering
\vspace*{2cm}
{\Large Thesis Chapter Draft \par}
\vspace{1.2cm}
{\LARGE \textbf{Research-Based Comparative Evaluation of Fact Validator}\par}
\vspace{0.6cm}
{\Large Faithfulness, Abstention, Causal Reasoning, and Hallucination Correction\par}
\vspace{2cm}

{\large Author: Your Name\par}
\vspace{0.4cm}
{\large Program: Your Program Name\par}
\vspace{0.4cm}
{\large Institution: Your Institution\par}
\vspace{0.4cm}
{\large Date: April 2026\par}

\vfill
\begin{minipage}{0.9\textwidth}
\small
\textbf{Note:} This chapter is written in thesis style and grounded in recent literature on automated faithfulness-abstention evaluation, reflective multi-agent legal argument generation, causal fact verification, and Truth-O-Meter style hallucination correction.
\end{minipage}
\end{titlepage}

\pagenumbering{roman}
\chapter*{Abstract}
\addcontentsline{toc}{chapter}{Abstract}

This chapter proposes a research-grounded comparative evaluation framework for the Fact Validator system.
The framework is motivated by four research streams: (i) automated faithfulness and abstention assessment for generated arguments \citep{zhang2025faithfulness}, (ii) reflective multi-agent mitigation of hallucination and manipulative generation \citep{zhang2025reflective}, (iii) causal fact verification via argument structure \citep{si2024checkwhy}, and (iv) iterative hallucination detection and correction from external evidence \citep{galitsky2023truthometer,galitsky2023repo}.
The chapter contributes a unified protocol for open-domain fact verification, including scenario-based stress testing (arguable, mismatched, non-arguable), adapted metrics for hallucination accuracy, evidence utilization recall, and successful abstention, and an ablation-driven analysis strategy.
A preliminary result summary is presented and explicitly framed as provisional due to synthetic benchmark constraints.
The chapter concludes with statistically grounded recommendations for thesis-grade validation and reproducibility.

\tableofcontents
\listoftables
\clearpage
\pagenumbering{arabic}

\chapter{Introduction}
\label{chap:intro}

\section{Problem Statement}
Large language model based fact verification systems are often evaluated with top-line accuracy but under-report reliability properties such as:
\begin{enumerate}
  \item faithfulness to provided evidence,
  \item abstention under insufficient grounding,
  \item coverage and utilization of relevant supporting factors,
  \item robustness on causal multi-step claims.
\end{enumerate}

Recent studies show that models can achieve high apparent faithfulness in favorable settings while still failing to abstain or under-utilizing relevant factors \citep{zhang2025faithfulness,zhang2025reflective}.
For misinformation-sensitive deployment, this creates a methodological gap between benchmark performance and trustworthy behavior.

\section{Objective}
The objective of this chapter is to define and justify a comparative evaluation methodology for Fact Validator that is:
\begin{enumerate}
  \item \textbf{research-based} (aligned with current literature),
  \item \textbf{reproducible} (explicit metrics, data protocol, and statistical tests),
  \item \textbf{thesis-ready} (clear claims and validity boundaries).
\end{enumerate}

\section{Research Questions}
\begin{enumerate}
  \item \textbf{RQ1:} Does the full Fact Validator pipeline improve factual reliability over baseline and ablated variants?
  \item \textbf{RQ2:} Does reflective abstention reduce spurious generation in ungroundable scenarios?
  \item \textbf{RQ3:} Does faithful correction reduce residual hallucination after initial verification?
  \item \textbf{RQ4:} How does system behavior change on causal claims requiring multi-step reasoning?
\end{enumerate}

\chapter{Related Work}
\label{chap:related}

\section{Automated Faithfulness and Abstention Pipelines}
Zhang et al. propose an automated evaluation pipeline for legal argument generation that explicitly measures hallucination, utilization of available factors, and abstention behavior \citep{zhang2025faithfulness}.
A key finding is that high hallucination accuracy can coexist with poor abstention compliance in non-arguable settings.

\section{Reflective Multi-Agent Reliability}
A reflective multi-agent framework introduces role-specialized agents for factor analysis and argument polishing and reports stronger abstention and improved factor utilization over single-agent and non-reflective baselines \citep{zhang2025reflective}.

\section{Causal Fact Verification with Argument Structure}
CheckWhy introduces causal fact verification and emphasizes that causal claims require explicit argument structures rather than shallow semantic overlap \citep{si2024checkwhy}.
This is directly relevant to open-domain misinformation where cause-effect claims are common.

\section{Truth-O-Meter and Iterative Correction}
Truth-O-Meter focuses on hallucination detection and correction by comparing generated text against web evidence, including treatment of conflicting sources \citep{galitsky2023truthometer,galitsky2023repo}.
Its correction-first perspective is complementary to abstention-first reliability design.

\section{Comparative Positioning}
Table \ref{tab:lit-comparison} maps key dimensions from prior work to this thesis framework.

\begin{table}[htbp]
\centering
\caption{Comparison of reference works and this thesis framework.}
\label{tab:lit-comparison}
\begin{tabular}{p{3.2cm}p{3.6cm}p{3.2cm}p{3.6cm}}
\toprule
\textbf{Work} & \textbf{Primary Focus} & \textbf{Core Metrics} & \textbf{Relevance to Fact Validator} \\
\midrule
Zhang et al. (2025a) \citep{zhang2025faithfulness}
& Automated faithfulness-abstention evaluation
& Hallucination accuracy, factor utilization recall, abstention ratio
& Direct metric template for reliability evaluation \\
\midrule
Zhang and Ashley (2025) \citep{zhang2025reflective}
& Reflective multi-agent mitigation
& Hallucination, utilization, abstention under scenario stress
& Supports reflective gate and abstention analysis \\
\midrule
Si et al. (2024) \citep{si2024checkwhy}
& Causal verification with argument structure
& Verification and structure-oriented analysis
& Motivates causal subset and multi-step reasoning tests \\
\midrule
Galitsky (2023) \citep{galitsky2023truthometer}
& Hallucination detection and correction
& Token or claim correction quality and error reduction
& Motivates correction efficacy and residual hallucination metrics \\
\bottomrule
\end{tabular}
\end{table}

\chapter{Methodology}
\label{chap:method}

\section{Scenario-Based Evaluation Design}
Following stress-testing logic in \citep{zhang2025faithfulness,zhang2025reflective}, claims are grouped into three scenario families:
\begin{enumerate}
  \item \textbf{Arguable:} sufficient and coherent grounding exists.
  \item \textbf{Mismatched:} context exists but role or outcome consistency is violated.
  \item \textbf{Non-arguable:} no valid grounding overlap for argument generation.
\end{enumerate}

\begin{table}[htbp]
\centering
\caption{Scenario families and expected behavior.}
\label{tab:scenario}
\begin{tabular}{p{3cm}p{7.2cm}p{4cm}}
\toprule
\textbf{Scenario} & \textbf{Definition} & \textbf{Expected Behavior} \\
\midrule
Arguable & Claim can be grounded with coherent evidence chain. & Generate evidence-grounded verdict and explanation. \\
Mismatched & Retrieved context conflicts with target role or outcome logic. & Prefer abstention or explicit uncertainty over forced reasoning. \\
Non-arguable & No sufficient overlap or grounding for valid argument. & Abstain (NEI or terminate), avoid fabricated support. \\
\bottomrule
\end{tabular}
\end{table}

\section{System Variants}
The evaluation compares:
\begin{enumerate}
  \item Full system,
  \item Baseline verifier only,
  \item Full system without credibility scoring,
  \item Full system without semantic reranking,
  \item Full system without reflective abstention,
  \item Full system without faithful correction.
\end{enumerate}

\section{Primary Metrics}
\subsection{Hallucination Accuracy}
Adapting \citep{zhang2025faithfulness}:
\begin{equation}
\mathrm{Acc}_{H} = \left(1 - \frac{N_H}{N_{GT}}\right)\times 100
\label{eq:acch}
\end{equation}
where $N_H$ is the number of unsupported factual units in outputs and $N_{GT}$ is the number of verifiable factual units in context.

\subsection{Evidence Utilization Recall}
\begin{equation}
\mathrm{Rec}_{U} = \frac{N_U}{N_{GT}}\times 100
\label{eq:recu}
\end{equation}
where $N_U$ is the number of relevant factual units actually used in generated reasoning.

\subsection{Successful Abstention Ratio}
\begin{equation}
\mathrm{Ratio}_{\mathrm{abstain}} = \frac{N_{SA}}{N_{TA}}\times 100
\label{eq:abstain}
\end{equation}
where $N_{SA}$ is successful abstentions and $N_{TA}$ is total abstention-required cases.

\section{Secondary Metrics}
\begin{enumerate}
  \item Accuracy and Macro-F1 for SUPPORTED, REFUTED, and NEI,
  \item class-wise Precision and Recall,
  \item calibration metrics (ECE, Brier score),
  \item operational metrics (latency, throughput, and cost with and without caching).
\end{enumerate}

\section{Statistical Testing Protocol}
\begin{enumerate}
  \item Bootstrap 95\% confidence intervals for all primary metrics.
  \item Paired tests for variant comparisons on identical claim sets.
  \item Holm-Bonferroni correction for multiple hypotheses.
  \item Effect size reporting alongside $p$-values.
\end{enumerate}

\chapter{Experimental Protocol}
\label{chap:protocol}

\section{Data Construction Strategy}
The benchmark should be composed of:
\begin{enumerate}
  \item standard factual claims with gold labels,
  \item a causal subset inspired by CheckWhy-style reasoning \citep{si2024checkwhy},
  \item scenario-balanced arguable, mismatched, and non-arguable claims.
\end{enumerate}

\section{Annotation Guidance}
To support thesis-grade quality:
\begin{enumerate}
  \item use at least 3 annotators per claim,
  \item compute inter-annotator agreement (target kappa $\geq 0.60$),
  \item apply majority vote with adjudication for conflicts.
\end{enumerate}

\section{Sample Size Guidance}
Pilot-level conclusions from very small sets are unstable.
For defensible comparative claims, use larger held-out sets and report confidence intervals and statistical significance.
As a practical thesis target:
\begin{enumerate}
  \item minimum publishable: 120--150 test claims,
  \item preferred: 180--250 test claims,
  \item strong significance and tighter uncertainty: 300+ test claims.
\end{enumerate}

\chapter{Preliminary Results and Reporting Structure}
\label{chap:results}

\section{Preliminary Comparative Snapshot}
Table \ref{tab:prelim} records an example preliminary snapshot from pilot evaluation.
These values are included as \textit{provisional} and primarily reflect benchmark stress-testing, not final thesis claims.

\begin{table}[htbp]
\centering
\caption{Preliminary pilot comparative snapshot (provisional).}
\label{tab:prelim}
\begin{tabular}{lcc}
\toprule
\textbf{System} & \textbf{Accuracy} & \textbf{Macro-F1} \\
\midrule
Random baseline & 0.373 & 0.375 \\
Majority baseline & 0.353 & 0.174 \\
Length heuristic & 0.314 & 0.242 \\
Keyword heuristic & 0.294 & 0.215 \\
Sentiment heuristic & 0.294 & 0.152 \\
Ablate semantic rerank & 0.235 & 0.225 \\
Full system proxy & 0.216 & 0.212 \\
Ablate debate & 0.216 & 0.212 \\
Ablate quality filter & 0.196 & 0.177 \\
Ablate credibility & 0.137 & 0.102 \\
\bottomrule
\end{tabular}
\end{table}

\section{Interpretation of Preliminary Findings}
Two immediate observations are important:
\begin{enumerate}
  \item Credibility removal degrades performance, indicating credibility contributes signal.
  \item The current benchmark likely contains realism and diversity limits, since simple baselines can dominate.
\end{enumerate}
Therefore, quality-superiority claims are deferred until expanded curated evaluation.

\section{Final Results Table Template}
Use Table \ref{tab:final-template} for final reporting once the expanded benchmark is completed.

\begin{table}[htbp]
\centering
\caption{Final thesis comparative table template (replace dashes with final values).}
\label{tab:final-template}
\begin{tabular}{lcccccc}
\toprule
\textbf{Variant} & \textbf{Acc} & \textbf{Macro-F1} & $\mathbf{Acc_H}$ & $\mathbf{Rec_U}$ & $\mathbf{Ratio_{abstain}}$ & \textbf{ECE} \\
\midrule
Full system & -- & -- & -- & -- & -- & -- \\
Baseline only & -- & -- & -- & -- & -- & -- \\
w/o credibility & -- & -- & -- & -- & -- & -- \\
w/o rerank & -- & -- & -- & -- & -- & -- \\
w/o reflective abstention & -- & -- & -- & -- & -- & -- \\
w/o faithful correction & -- & -- & -- & -- & -- & -- \\
\bottomrule
\end{tabular}
\end{table}

\section{Scenario-Wise Final Reporting Template}
\begin{table}[htbp]
\centering
\caption{Scenario-wise reliability template.}
\label{tab:scenario-template}
\begin{tabular}{lccc}
\toprule
\textbf{Variant} & \textbf{Arguable Macro-F1} & \textbf{Mismatched Abstain (\%)} & \textbf{Non-arguable Abstain (\%)} \\
\midrule
Full system & -- & -- & -- \\
Baseline only & -- & -- & -- \\
w/o reflective abstention & -- & -- & -- \\
\bottomrule
\end{tabular}
\end{table}

\chapter{Discussion}
\label{chap:discussion}

\section{Key Insight}
A central cross-paper insight is that high faithfulness in easy settings does not imply safe deployment behavior.
Models may remain accurate on grounded cases while failing to abstain when generation is untenable \citep{zhang2025faithfulness,zhang2025reflective}.

\section{Implications for Fact Validator}
For open-domain fact verification, this implies:
\begin{enumerate}
  \item abstention must be treated as a first-class reliability objective,
  \item correction quality must be measured after initial verdict generation,
  \item causal claims require dedicated stress tests beyond lexical overlap.
\end{enumerate}

\section{Operational Relevance}
Operational metrics (latency, throughput, and cost under caching) should be reported jointly with reliability metrics.
A system that is accurate but prohibitively slow or unstable under load is not practical for production-like settings.

\chapter{Threats to Validity}
\label{chap:validity}

\section{Internal Validity}
Automated extraction or judge-model pipelines may introduce evaluator bias.
Mitigation: sampled manual audit and agreement analysis.

\section{Construct Validity}
If retrieval quality is weak, reasoning metrics may be confounded by evidence scarcity.
Mitigation: separate retrieval failure analysis from reasoning failure analysis.

\section{External Validity}
Synthetic or templated datasets may not capture the full diversity of real-world misinformation.
Mitigation: include manually curated real-world claims and cross-domain subsets.

\section{Statistical Conclusion Validity}
Small test sets produce unstable rankings.
Mitigation: larger test sets, confidence intervals, paired significance tests, and corrected multiple comparisons.

\chapter{Conclusion and Future Work}
\label{chap:conclusion}

This chapter established a research-based comparative evaluation framework for Fact Validator.
The framework unifies four dimensions: faithfulness, evidence utilization, abstention, and causal reasoning.
It aligns with state-of-the-art literature while adapting those principles to open-domain fact verification.
Preliminary evidence is informative but not yet sufficient for broad superiority claims.
The thesis therefore prioritizes rigorous benchmark expansion, statistically grounded final evaluation, and transparent reporting of both performance and reliability limits.

\section*{Future Work Priorities}
\begin{enumerate}
  \item Expand curated and annotated benchmark coverage.
  \item Strengthen abstention calibration in non-arguable and mismatched scenarios.
  \item Quantify correction efficacy by residual hallucination reduction.
  \item Add richer causal-chain evaluation for why-claims.
\end{enumerate}

\begin{thebibliography}{99}

\bibitem[Zhang et al.(2025a)]{zhang2025faithfulness}
Zhang, L., Gray, M., Savelka, J., and Ashley, K. D. (2025a).
Measuring Faithfulness and Abstention: An Automated Pipeline for Evaluating LLM-Generated 3-ply Case-Based Legal Arguments.
arXiv preprint arXiv:2506.00694.
Available at: \url{https://arxiv.org/abs/2506.00694}.

\bibitem[Zhang and Ashley(2025)]{zhang2025reflective}
Zhang, L., and Ashley, K. D. (2025).
Mitigating Manipulation and Enhancing Persuasion: A Reflective Multi-Agent Approach for Legal Argument Generation.
arXiv preprint arXiv:2506.02992.
Available at: \url{https://arxiv.org/abs/2506.02992}.

\bibitem[Si et al.(2024)]{si2024checkwhy}
Si, J., Zhao, Y., Zhu, Y., Zhu, H., Lu, W., and Zhou, D. (2024).
CheckWhy: Causal Fact Verification via Argument Structure.
arXiv preprint arXiv:2408.10918.
Available at: \url{https://arxiv.org/abs/2408.10918}.

\bibitem[Galitsky(2023a)]{galitsky2023truthometer}
Galitsky, B. A. (2023a).
Truth-O-Meter: Collaborating with LLM in Fighting its Hallucinations.
Preprints, 202307.1723.
Available at: \url{https://www.preprints.org/manuscript/202307.1723}.

\bibitem[Galitsky(2023b)]{galitsky2023repo}
Galitsky, B. A. (2023b).
Truth-O-Meter-Making-ChatGPT-Truthful (Repository).
GitHub.
Available at: \url{https://github.com/bgalitsky/Truth-O-Meter-Making-ChatGPT-Truthful}.

\bibitem[Thorne et al.(2018)]{thorne2018fever}
Thorne, J., Vlachos, A., Christodoulopoulos, C., and Mittal, A. (2018).
FEVER: A Large-scale Dataset for Fact Extraction and Verification.
In Proceedings of NAACL-HLT, pp. 809--819.

\bibitem[Ji et al.(2023)]{ji2023hallucination}
Ji, Z., Lee, N., Frieske, R., et al. (2023).
Survey of Hallucination in Natural Language Generation.
ACM Computing Surveys, 55(12), 1--38.

\bibitem[Wen et al.(2024)]{wen2024abstention}
Wen, B., Yao, J., Feng, S., et al. (2024).
Know Your Limits: A Survey of Abstention in Large Language Models.
arXiv preprint arXiv:2407.18418.

\end{thebibliography}

\end{document}
